\begin{example}
\label{example-spec-kxy}
In this example we describe $X = \Spec(k[x, y])$
when $k$ is an arbitrary field.
Clearly $(0)$ is prime, and any principal ideal generated by an
irreducible polynomial will also be a prime since $k[x, y]$ is a
unique factorization domain. Now assume $\mathfrak p$ is an
element of $X$ that is not principal. Since $k[x, y]$ is a
Noetherian UFD, the prime ideal $\mathfrak p$ can be generated
by a finite number of irreducible polynomials $(f_1, \ldots, f_n)$.
Now, I claim that if $f, g$ are irreducible polynomials in $k[x, y]$
that are not associates, then $(f, g) \cap k[x] \neq 0$. To do this,
if either $f$ or $g$ lies in $k[x]$, that nonzero polynomial
itself belongs to $(f,g)\cap k[x]$ and proves the assertion.
Otherwise both have positive degree in $y$. Each is primitive as
a polynomial in $y$ over $k[x]$, since a nonunit content would
factor the original irreducible polynomial in $k[x,y]$.
Gauss' lemma therefore makes both irreducible in $k(x)[y]$.
In this case it is enough to show that $f$ and $g$ are relatively prime when
viewed in $k(x)[y]$. In this case, $k(x)[y]$ is a Euclidean domain,
so by applying the Euclidean algorithm and clearing denominators, we
obtain $p = af + bg$ for $0 \ne p \in k[x]$ and $a, b \in k[x, y]$. Thus, assume this is not
the case, that is, that some nonunit $h \in k(x)[y]$ divides both
$f$ and $g$. Then, by Gauss's lemma, for some $a, b \in k(x)$ we
have $ah | f$ and $bh | g$ for $ah, bh \in k[x, y]$. By
irreducibility, $ah = f$ and
$bh = g$ (since $h \notin k(x)$). So, back in $k(x)[y]$, $f, g $
are associates, as $\frac{a}{b} g = f$. Since
$k(x)$ is the fraction field of $k[x]$, we can write $g = \frac{r}{s} f $
for elements $r , s \in k[x]$ sharing no common factors. This
implies that $sg = rf$ in $k[x, y]$ and so $s$ must divide $f$
since $k[x,y]$ is a UFD. If $s$ is a unit in $k[x]$, write
$s=c\in k^\times$. Then $g=(r/c)f$; irreducibility of $g$
forces $r/c$ to be a unit, contradicting that $f$ and $g$ are
not associates. If $s$ is not a unit, the divisibility $s\mid f$
and irreducibility of $f$ give $f=us$ for a unit $u\in k^\times$.
This would put $f$ in $k[x]$, contrary to its positive $y$-degree.
Both alternatives contradict the existence of the common divisor.
Thus $f$ and $g$ are relatively prime in $k(x)[y]$, and the
cleared Bezout identity gives a nonzero polynomial in $(f,g)\cap k[x]$.
To distinguish it from a selected irreducible factor below, denote
this resulting polynomial by $P$.
The irreducible generators of the nonprincipal prime $\mathfrak p$
include two which are not associates, since otherwise they generate
a principal ideal. Applying the argument to those two puts
$P\in\mathfrak p\cap k[x]$. It cannot be a nonzero constant,
because $\mathfrak p$ is proper. Write its full factorization as
$P=c\prod_{j=1}^d p_j$, with $c\in k^\times$ and irreducible
$p_j\in k[x]$, repeating factors according to their multiplicities.
Primality gives some $p_j\in\mathfrak p$; denote this chosen
factor by $p$. Thus $\mathfrak p$ contains the irreducible
polynomial $p\in k[x]\subset k[x,y]$.
By the above, $\mathfrak p$ corresponds to a prime in the ring
$k[x,y]/(p)\cong(k[x]/(p))[y]=k(\alpha)[y]$, where
$\alpha=x+(p)$ in the field $k[x]/(p)$.
The isomorphism sends $F(x,y)+(p)$ to $F(\alpha,y)$:
coefficientwise reduction is surjective and its kernel consists
exactly of the polynomials whose coefficients are multiples of $p$.
Every element of $k[x]/(p)$ is a polynomial in $\alpha$, and this
quotient is a field, so it is $k(\alpha)$. No leading coefficient
of the chosen $p$ is changed or omitted. This is a
PID, and so any prime ideal corresponds to $(0)$ or an
irreducible polynomial in $k(\alpha)[y]$. Thus, $\mathfrak p$
is of the form $(p)$ or $(p, f)$ where $f$ is a
polynomial in $k[x, y]$ that is irreducible in the quotient
$k[x, y]/(p)$.
\end{example}